\section{特征与 Frobenius 自同态}\label{sec:galois-2-4}
\subsection{域的特征}\label{sec:galois-2-4-1}
设 $A$ 为交换环. 存在唯一的环态射 $\mathbb Z\longrightarrow A$. 事实上，这样的态射必须将 $1$ 映为 $1_A$，将 $n$ 映为 $1_A+\cdots+1_A$（$n$ 项）. 此态射的核称为 $A$ 的\newterm{galois-characteristic-ideal}{特征理想}{characteristic ideal}. 当 $A=k$ 为域时，可能出现以下两种情形：
\begin{itemize}
\item 特征理想为零，此时称 $k$ 的\newterm{galois-characteristic}{特征}{characteristic}为零.
\item 特征理想为素理想，因而由唯一的素数 $p$ 生成，此时称 $k$ 的特征为 $p$.
\end{itemize}
\begin{remark*}
若 $p$ 为素数，更一般地，当环 $A$ 的特征理想等于 $(p)$ 时，称 $A$ ``特征为 $p$''. 此时态射 $\mathbb Z\longrightarrow A$ 经由 $\mathbb F_p:=\mathbb Z/p\mathbb Z$ 分解，其中 $\mathbb F_p$ 是含 $p$ 个元素的有限域，因而 $A$ 是 $\mathbb F_p$-代数. 反之，每个 $\mathbb F_p$-代数都是特征为 $p$ 的环.
\end{remark*}
\subsection{素子域}\label{sec:galois-2-4-2}
域 $k$ 的\newterm{galois-prime-subfield}{素子域}{prime subfield}是指 $k$ 的最小子域，即 $k$ 的所有子域的交. 可能出现以下两种情形：
\begin{itemize}
\item 若 $k$ 的特征为零，则 $k$ 包含 $\mathbb Z$，因而包含 $\operatorname{Frac}(\mathbb Z)=\mathbb Q$，故 $k$ 的素子域为 $\mathbb Q$.
\item 若 $k$ 的特征为 $p>0$，则 $k$ 包含 $\mathbb F_p$，它就是 $k$ 的素子域.
\end{itemize}
\subsection{Frobenius 自同态}\label{sec:galois-2-4-3}
\begin{proposition*}
设 $A$ 是特征为 $p$ 的环. 则映射
\begin{align*}
F_A:A\longrightarrow A,\qquad a\longmapsto a^p
\end{align*}
是 $\mathbb F_p$-代数自同态，称为 $A$ 的\newterm{galois-frobenius-endomorphism}{Frobenius 自同态}{Frobenius endomorphism}.
\end{proposition*}
\begin{proof}
设 $a,b\in A$. 显然 $(ab)^p=a^pb^p$. 另一方面，
\begin{align*}
(a+b)^p=\sum_{k=0}^p\binom pk a^{p-k}b^k,\qquad \binom pk=\frac{p!}{k!(p-k)!}.
\end{align*}
当 $0<k<p$ 时，$p$ 既不整除 $k!$，也不整除 $(p-k)!$. 但 $p$ 整除 $p!$，因而整除 $\binom pk$. 所以在 $A$ 中有 $(a+b)^p=a^p+b^p$.
\end{proof}
\begin{remark*}
Frobenius 自同态与任意 $\mathbb F_p$-代数态射都``交换''：若 $\varphi:A\longrightarrow B$ 是这样的态射，则 $\varphi\circ F_A=F_B\circ\varphi$.
\end{remark*}
对于特征为 $p$ 的域 $k$，记
\begin{align*}
k^F:=\{x\in k,F_k(x)=x\}
\end{align*}
为 Frobenius 自同态的不动点集. 由于 $F_k$ 是域自同态，它是 $k$ 的子域.
\begin{lemma*}
$k^F$ 就是 $k$ 的素子域 $\mathbb F_p$.
\end{lemma*}
\begin{proof}
对于 $x\in k$，在 $k[X]$ 中有
\begin{align*}
F_k(x)=x\ \Longleftrightarrow\ x^p=x\ \Longleftrightarrow\ (X-x)\mid(X^p-X).
\end{align*}
另一方面，对每个 $a\in\mathbb F_p$ 都有 $a^p=a$. 当 $a$ 遍历 $\mathbb F_p$ 时，不可约多项式 $X-a$ 两两互素，故在 $\mathbb F_p[X]$ 及 $k[X]$ 中都有分解 $X^p-X=\prod_{a\in\mathbb F_p}(X-a)$. 因而 $x\in\mathbb F_p$.
\end{proof}
更一般地，对每个 $r\in\mathbb N^*$，自同态 $F_k^r=F_k\circ F_k\circ\cdots\circ F_k$ 的不动点集
\begin{align*}
k^{F^r}:=\{x\in k,F_k^r(x)=x\}
\end{align*}
是 $k$ 的子域. 当 $k$ 是 $\mathbb F_p$ 的代数闭包时，这一情形尤其值得关注.
\subsection{有限域}\label{sec:galois-2-4-4}
选取 $\mathbb F_p$ 的一个代数闭包 $\overline{\mathbb F}_p$，记其 Frobenius 自同构为 $F$.
\begin{theorem*}
域 $\overline{\mathbb F}_p^{F^r}$ 是多项式 $X^{p^r}-X$ 在 $\mathbb F_p$ 上的分裂域. 反之，$\mathbb F_p$ 的每个有限扩张 $k$ 都是多项式 $X^{p^{[k:\mathbb F_p]}}-X$ 在 $\mathbb F_p$ 上的分裂域.
\end{theorem*}
\begin{proof}
对 $x\in\overline{\mathbb F}_p$，有 $F^r(x)=x$ 当且仅当 $x$ 是 $X^{p^r}-X$ 的根. 因此 $\overline{\mathbb F}_p^{F^r}$ 是 $X^{p^r}-X$ 在 $\overline{\mathbb F}_p$ 中的根的集合. 由于它是域，特别地，它就是 $X^{p^r}-X$ 的分裂域.

反之，设 $k$ 是 $\mathbb F_p$ 的有限扩张，记其在 $\mathbb F_p$ 上的维数为 $r:=[k:\mathbb F_p]$. 则 $k$ 为有限集，基数为 $|k|=p^r$，故其乘法群 $k^\times$ 的阶为 $p^r-1$，从而每个 $x\in k^\times$ 都满足 $x^{p^r-1}=1$. 因此 $k$ 的每个元素 $x$ 都是多项式 $X(X^{p^r-1}-1)=X^{p^r}-X$ 的根. 特别地，$k$ 是此多项式的分裂域.
\end{proof}
由于每个有限域都是其素子域的有限扩张，这一定理给出了``构造''全部有限域的方法. 为完善这一定理，还需计算 $\overline{\mathbb F}_p^{F^r}$ 的基数，这等价于数出 $X^{p^r}-X$ 的根（至多有 $p^r$ 个）. 下一节将完成此事.

